\section{结论与政策建议}

本文基于 BACI HS07 2008--2024 年中国进口贸易数据，研究四类半导体相关产品的进口来源结构演化，并构建 SIRI 风险指数评价供应链风险。主要结论如下。

第一，中国半导体相关产品进口来源结构在 2018 年后出现明显重组。848620、854231 和 854232 的美国进口份额下降，854239 则小幅上升，说明不能将所有半导体相关产品概括为同一种变化方向。

第二，不同产品的风险暴露存在显著异质性。设备类产品的替代来源更多集中于日本、荷兰、新加坡等设备强国；集成电路类产品则更多体现亚洲生产网络内部调整。SIRI 指数显示，2024 年 854239 和 854232 的综合风险相对更高。

第三，固定效应回归结果不支持将来源结构变化简单归因为美国出口管制的严格因果效应。更稳妥的表述是：政策冲击背景下，中国半导体相关产品进口来源结构发生重组，美国份额在部分产品中承压，但绝对进口额变化仍受到产业周期、需求扩张和第三方产能变化等因素共同影响。

基于以上发现，本文提出三点建议。其一，对设备类和集成电路类产品采取差异化风险监测，不宜用单一产品结论推断整体半导体供应链。其二，持续跟踪来源集中度和替代来源结构，尤其关注高 SIRI 产品的主要来源国集中风险。其三，将风险指数作为动态监测工具，与企业层面供应链信息和更高频贸易数据结合，提高关键产品进口安全预警能力。
